\documentclass[11pt]{article}
\usepackage{amsmath,amssymb,amsthm}
\newtheorem{lemma}{Lemma}
\title{Shared high-order cancellations and their unresolved readout problem}
\author{}
\date{}
\begin{document}
\maketitle

\paragraph{Verification and scope.}
The two identities below have elementary complete written proofs. The
companion checker verifies finite rational instances. Neither identity
is a superconstant compiler. In particular, a large common order of
vanishing is not a count of independently selectable signs.

\begin{lemma}[Barycentric shared orders]
Let $x_1,\ldots,x_H$ be distinct real numbers and put
$a_h=\prod_{g\ne h}(x_h-x_g)^{-1}$. For any real $u$ and integer
$k\geq1$ with $2k-1<H-1$,
\[
 \sum_{h=1}^H a_h(1+ux_h)^k x_h^{k-1}=0.
\]
Consequently the normalized reverse query
\[
 F_{u,t}(\delta)=\sum_h
 \frac{\delta a_h}{t/(1+ux_h)-\delta x_h}
\]
has all its coefficients of orders $1\leq k<H/2$ equal to zero,
simultaneously for every $u$ for which all $1+ux_h$ are positive and
every positive $t$. Its denominators are positive for sufficiently
small positive $\delta$.
\end{lemma}
\begin{proof}
Lagrange interpolation implies $\sum_h a_h Q(x_h)=0$ for every
polynomial $Q$ of degree below $H-1$: compare coefficients of $X^{H-1}$
in its interpolation formula. Apply this to
$Q(X)=(1+uX)^kX^{k-1}$, of degree at most $2k-1$.
The coefficient of $\delta^k$ in the displayed reciprocal sum is
$t^{-k}\sum_h a_h(1+ux_h)^kx_h^{k-1}$, by the geometric series.
\end{proof}

For varying first directions $(A,P)$, uniform cancellation through
order $D-1$ in this family is equivalent to the mixed moments
\[
 \sum_h A_h x_h^l P_h^{k-1}=0
 \quad(1\leq k<D,\ 0\leq l\leq k),
\]
with the usual omission of vacuous or repeated conditions. This uses
only $O(D^2)$ polynomial conditions, irrespective of the number of
query values $u$. However, the next coefficient remains the single
degree-$D$ polynomial
\[
 R_D(u)=\sum_{l=0}^D\binom Dl u^l
            \sum_h A_hx_h^lP_h^{D-1}.
\]
Thus this leading coefficient cannot supply arbitrary signs at more
than $D+1$ ordered $u$ values: the fully alternating pattern would
require more than $D$ distinct real zeros. This is only a limitation
of the displayed leading readout, not of all perturbations of the
original rational family.

Query-specific numerator offsets can vanish on the fixed barycentric
curve to high order. For a query with reciprocal vector $w_h>0$, its
offset coefficients $b_h$ must satisfy
\[
 \sum_h b_hw_h^{k+1}x_h^k=0\qquad(0\leq k<D).
\]
These are at most $D$ linear conditions in its $H$ freely chosen
offsets. They leave many offset choices for a fixed curve. What is
not proved is that these cancellations persist on a sufficiently rich
family of variable first directions while their higher terms provide
independent signs at a superlinear number of queries. Imposing them
separately for many queries can reintroduce the original number of
moment constraints. No independence or dimension conclusion is inferred
merely from counting the coefficients in one query.

\begin{lemma}[Walsh shared orders]
Let $H=2^d$, index heads by $h\in\{-1,1\}^d$, and put
$P_h=\sum_{j=1}^d x_jh_j$. For $S\subseteq[d]$ let
$\chi_S(h)=\prod_{j\in S}h_j$. Then
\[
 \sum_h\chi_S(h)P_h^k=0\quad(k<|S|),\qquad
 \sum_h\chi_S(h)P_h^{|S|}
 =2^d|S|!\prod_{j\in S}x_j.
\]
\end{lemma}
\begin{proof}
Expand the power by the multinomial theorem. Summing over each sign
$h_j$ annihilates a monomial unless its exponent is odd for $j\in S$
and even for $j\notin S$. This requires degree at least $|S|$. At
that degree, every coordinate in $S$ appears exactly once and every
other coordinate has exponent zero, giving the stated coefficient.
\end{proof}

Approximately half the Walsh masks have a common vanishing order
proportional to $\log H$. Their displayed leading signs are nevertheless
products of only $d$ scalar signs. This identity therefore demonstrates
shared cancellation, not independent readout. Higher perturbations of
the poles or residues require a separate universality argument.

\paragraph{Exact finite progress.}
The nonlinear cone in
\texttt{reverse-unbalanced-h4-cubic.tex} does yield a complete finite
improvement: four heads realize every ten-query sign pattern with a
strict whole-cube embedding for arbitrarily many rows. The proof uses
the cubic remainder after two cancellations and exact rational witnesses
on a quadratic moment cone. It does not extrapolate either of the
shared-order identities above into an asymptotic construction.
\end{document}
